% Unidad: definición de espacio normado (castellano).
%
% Sin preámbulo. La unidad no sabe a qué salida va: los interruptores deciden.
% Compárese con la versión heredada, que arrastraba su propio \documentclass,
% su propio babel (equivocado) y cuatro \import con rutas relativas.

\begin{frame}
\didactatitle{Definición y propiedades de los espacios normados}

La métrica usual en $\mathbb{R}$ es la distancia euclídea entre dos puntos
$x, y \in \mathbb{R}$, dada por
\[
  d(x,y) = |x-y|.
\]
Este concepto se generaliza a espacios vectoriales mediante la noción de norma.

\onlynotes{%
  La idea detrás del concepto de norma es generalizar las propiedades más
  útiles de la distancia euclídea en $\mathbb{R}$: que mide algo no negativo,
  que escala con el factor de escala, y que respeta la desigualdad triangular.
}

\begin{definition}[Espacio normado]
Un \keyterm{espacio normado} es un par $(E, \|\cdot\|)$, donde $E$ es un
espacio vectorial sobre el cuerpo $\mathbb{K}$ ($\mathbb{R}$ o $\mathbb{C}$) y
$\|\cdot\|$ es una función llamada \keyterm{norma} que satisface:
\begin{itemize}
  \item \textbf{Positividad:} $\|x\| \geq 0$, y $\|x\| = 0 \iff x = 0$.
  \item \textbf{Homogeneidad:} $\|\lambda x\| = |\lambda|\,\|x\|$ para todo
        $\lambda \in \mathbb{K}$ y $x \in E$.
  \item \textbf{Desigualdad triangular:} $\|x+y\| \leq \|x\| + \|y\|$ para
        todo $x, y \in E$.
\end{itemize}
Cuando no haya ambigüedad denotaremos el espacio normado simplemente por $E$.
\end{definition}

\begin{teaching}[Ritmo]
Quince minutos. Merece la pena detenerse en la homogeneidad: casi nadie ve por
qué hace falta el valor absoluto en $|\lambda|$ hasta que se les pide
$\lambda = -1$.
\end{teaching}
\end{frame}

\begin{frame}
\didactatitle{Ejemplos de espacios normados}

Un espacio normado permite medir la \keyterm{longitud} de los vectores y
definir conceptos como \keyterm{distancia} y \keyterm{convergencia}.

\dpause

\begin{example}
En $\mathbb{R}^n$ podemos definir distintas normas. Las más comunes son
\[
  \|x\|_1 = \sum_{i=1}^n |x_i|,
  \qquad
  \|x\|_2 = \Bigl(\sum_{i=1}^n |x_i|^2\Bigr)^{1/2},
  \qquad
  \|x\|_\infty = \max_{1 \leq i \leq n} |x_i|.
\]

\onlynotes{%
  \begin{itemize}
    \item $\|x\|_2$ es la \keyterm{norma euclídea}, y su métrica asociada es la
          distancia euclídea.
    \item $\|x\|_1$ se usa en optimización, por su simplicidad computacional.
    \item $\|x\|_\infty$ es útil en análisis numérico, cuando interesa
          controlar el \keyterm{valor máximo} de las componentes.
  \end{itemize}
}

También podemos definir normas en espacios más sofisticados, como el de las
funciones continuas en $[a,b]$ con
\begin{keyformula}
  \|f\| = \max_{x \in [a,b]} |f(x)|.
\end{keyformula}
\end{example}

\begin{commonmistake}
Escriben $\|\lambda x\| = \lambda \|x\|$ y pierden el valor absoluto. Sale
sistemáticamente en el primer parcial.
\end{commonmistake}
\end{frame}
